\documentclass[aspectratio=169]{beamer}

\usepackage[utf8]{inputenc}
\renewcommand{\familydefault}{\sfdefault}
\usepackage{amsmath}

\title{Functional Programming for rubyists}
\author{Bianca Gibson}
\institute{Ruby Conf AU 2016}
\date{}

\begin{document}

\frame{\titlepage}

\begin{frame}
  \begin{center}
    \Huge Rampant side effects
  \end{center}
\end{frame}

\begin{frame}
  \begin{center}
    \Huge Shifting sands of mutable state beneath your feet
  \end{center}
\end{frame}

\begin{frame}
  \begin{center}
    \Huge Enter Functional Programming
  \end{center}
\end{frame}

\begin{frame}
  \begin{figure}[p]
    \includegraphics[width=0.4\textwidth]{./assets/portrait.jpg}
  \end{figure}
\end{frame}

\begin{frame}
  \begin{center}
    \Huge What should you get out of this?
  \end{center}
  %\begin{itemize}
  %  \item What is the core of FP?
  %  \item How can it help with problems you've experienced?
  %  \item How to use it in ruby projects you're already working on
  %  \item It's not that hard and complicated
  %  \item No neckbeard required
  %\end{itemize}
\end{frame}

\begin{frame}
  \begin{center}
    \Huge What will we cover
  \end{center}
  %\begin{itemize}
  %  \item Immutability
  %  \item Putting off side effects
  %  \item Referential transparency
  %  \item Common traps when explaining to others and avoiding them
  %  \item How to learn more
  %\end{itemize}
\end{frame}

\begin{frame}
  \begin{center}
    \Huge Stop things changing behind my back
  \end{center}
\end{frame}
% tell gradle story here

\begin{frame}
  \begin{center}
    \Huge Immutability
  \end{center}
  %\begin{itemize}
  %  \item Don't change any values, use only constants
  %  \item Instead of changing create a new one
  %  \item Recursion and accumulators over loops and mutation
  %  \item Confusion over passing by value or reference? Gone.
  %  \item Ever thought about using a ! or not? That's worrying about mutation
  %\end{itemize}
\end{frame}

\begin{frame}
  \begin{figure}[p]
    \includegraphics[width=0.8\textwidth]{./assets/confused.jpg}
  \end{figure}
  \tiny Image from CollegeDegrees360, license https://creativecommons.org/licenses/by-sa/2.0/
\end{frame}

\begin{frame}
  \begin{center}
    \Huge Example
  \end{center}
\end{frame}

\begin{frame}
  \begin{center}
    \Huge What are side effects?
  \end{center}
  %\begin{itemize}
  %  \item When your code affects the outside world
  %  \item Printing
  %  \item Writing to disk
  %  \item Network calls
  %\end{itemize}
\end{frame}

\begin{frame}
  \begin{center}
    \Huge What do they have to do with FP?
  \end{center}
  %\begin{itemize}
  %  \item Separate your logic from your side effects
  %  \item Push your side effects out in to the edge of your system
  %  \item In the middle of your system make representations of the side effects
  %  \item Makes the core easier to test and reason about
  %  \item No side effects to stub $\rightarrow$ simpler tests
  %  \item Push hard testing out to the edge
  %\end{itemize}
\end{frame}

\begin{frame}
  \begin{center}
    \Huge Example
  \end{center}
\end{frame}

\begin{frame}
  \begin{center}
    \Huge Referential transparency
  \end{center}
  %\begin{itemize}
  %  \item The combination of immutability and no side effects
  %  \item When you call a function again you get the same result
  %  \item Includes no side effects
  %  \item You can substitute the saved result of the function for the call
  %  \item Everything a function does is represented by the result
  %  \item Nothing unexpected!
  %\end{itemize}
\end{frame}

\begin{frame}
  \begin{center}
    \Huge Example
  \end{center}
\end{frame}

\begin{frame}
  \begin{center}
    \Huge Teaching FP
  \end{center}
  %\begin{itemize}
  %  \item Don't get too mathematical
  %  \item Relate to stuff they already know
  %  \item Focus on examples
  %  \item Avoid the scary words like monad until they understand it
  %  \item Focus on advantages to them
  %  \item Pick your battles
  %\end{itemize}
\end{frame}

\begin{frame}
  \frametitle{Learning more}
  \begin{itemize}
    \item Functional Programming Principles in Scala on coursera
    \item Functional programming in Scala by Paul Chiusano and Rúnar Bjarnason
    \item Refactoring Ruby with Monads by Tom Stuart
    \item github.com/cwmyers/FunctionalTraining
    \item github.com/NICTA/course
  \end{itemize}
\end{frame}

\begin{frame}
  \begin{center}
    \Huge Questions?
  \end{center}
\end{frame}

\begin{frame}
  \begin{center}
    \Huge Stop things changing behind our backs
  \end{center}
\end{frame}

\begin{frame}
  \begin{center}
    \Huge Follow up and feedback to bianca.gibson@rea-group.com
  \end{center}
\end{frame}

\end{document}
